%====================================================================
%  Directional Congestion and Tidal Lanes
%
%  A self-contained statement of the directional quantitative spatial
%  commuting model: setup, equilibrium definition, master inventory of
%  variables/parameters/data sources, and the estimation pipeline
%  (lambda from a directed speed-flow regression; inversion of
%  fundamentals; nested-fixed-point exact-hat solver for counterfactuals).
%====================================================================
\documentclass[12pt]{article}

\usepackage[margin=1in]{geometry}
\usepackage{amsmath,amssymb,amsthm,mathtools}
\usepackage{graphicx,booktabs,tabularx,array,setspace,float,enumitem}
\usepackage{longtable}
\usepackage{natbib}
\usepackage[colorlinks=true,linkcolor=blue,citecolor=blue,urlcolor=blue]{hyperref}
\setstretch{1.15}
\setlength{\parindent}{1.5em}

\newtheorem{proposition}{Proposition}
\newtheorem{lemma}{Lemma}
\newtheorem{definition}{Definition}
\newtheorem{assumption}{Assumption}
\newtheorem{remark}{Remark}

\DeclareMathOperator*{\argmax}{arg\,max}

\newcommand{\N}{\mathcal N}
\newcommand{\E}{\mathcal E}
\newcommand{\R}{\mathcal R}
\newcommand{\Lbar}{\bar L}
\newcommand{\Wbar}{\bar W}
\newcommand{\ubar}{\bar u}
\newcommand{\Abar}{\bar A}
\newcommand{\tbar}{\bar t}
\newcommand{\tkl}{t_{kl}}
\newcommand{\tijk}{\tau_{ij}}
\newcommand{\obs}{\mathrm{obs}}
\newcommand{\cf}{\mathrm{cf}}

\title{Directional Congestion and Tidal Lanes:\\
Model Part}
\author{Yige Duan \and Xiaoruo Feng \and Yizhen Gu}
\date{\today}

\begin{document}
\maketitle

\begin{abstract}
\noindent
This paper develops a quantitative spatial commuting model in which every primitive of the road network is directed. Lane capacities, free-flow speeds, link travel times, and bilateral commuting cost indices are allowed to differ between $(k,l)$ and $(l,k)$; workers choose a residence, a workplace, and a route under Fréchet preferences; and equilibrium delivers gravity in commuting flows, gravity in directed link traffic, and a closed-form welfare scale. Three features distinguish this paper. The first is directional speed asymmetry at the segment level: observed peak-period imbalances between the two directions of a corridor enter the cost structure directly, rather than being subsumed in a symmetric travel-time index. The second is the policy experiment: tidal-lane reallocation reassigns lane capacity between directions within existing road width, requiring neither new land supply nor substantial physical construction, making it among the lowest-cost levers available to urban congestion management. The third is the data input: granular, high-frequency segment-level speed observations identify the direction-specific congestion wedge on each corridor and supply the directed cost primitives the model requires.
\end{abstract}

\tableofcontents
\clearpage

%====================================================================
\section{Setup}\label{sec:setup}
%====================================================================

\subsection{Locations and directed network}

The city is partitioned into $N$ spatial cells $\N=\{1,\dots,N\}$,
each indexed by its centroid.  The road network is a directed graph
$(\N,\E)$ with edge set $\E\subseteq\N\times\N$.  A directed edge
$(k,l)\in\E$ records an admissible move from cell $k$ to cell $l$.
$(k,l)\in\E$ does not imply $(l,k)\in\E$, and even when both edges
exist the two directions are distinct objects with distinct
primitives.

\subsection{Directed link primitives}

Each directed link $(k,l)\in\E$ carries three primitives:
\begin{itemize}[leftmargin=2em,itemsep=0.2em]
  \item physical length $\ell_{kl}>0$ (centroid-to-centroid);
  \item directed lane capacity $n_{kl}>0$;
  \item free-flow speed $v^{0}_{kl}>0$, with associated free-flow
        travel time $\tbar_{kl}\equiv\ell_{kl}/v^{0}_{kl}$.
\end{itemize}
Congested directed travel time $\tkl$ is endogenous and
determined by the congestion technology in
\S\ref{sec:objects}.  Directionality is permitted in every primitive:
$\ell_{kl}\ne\ell_{lk}$, $n_{kl}\ne n_{lk}$, $v^{0}_{kl}\ne v^{0}_{lk}$,
$\tkl\ne t_{lk}$ are all admitted.

\subsection{Iceberg cost convention}\label{ssec:iceberg}

The model uses a dimensionless iceberg link cost $\tkl\ge 1$ with the
\citet{AllenArkolakis2022} micro-foundation
\begin{equation}
  \tkl\;=\;\bigl(\ell_{kl}/v_{kl}\bigr)^{1/\theta},
  \qquad
  v_{kl}\;=\;v^{0}_{kl}\bigl(n_{kl}/\Xi_{kl}\bigr)^{\lambda},
  \label{eq:iceberg}
\end{equation}
where $v_{kl}$ is the realised (congested) directed speed and
$\Xi_{kl}$ is endogenous directed link traffic.  The iceberg cost
$\tkl$ enters utility multiplicatively as
$\prod_m t_{r_{m-1},r_m}$ and the kernel of the route system is
$\tkl^{-\theta}$.  The $\tkl$ here is scale invariance, rather than physical minutes. Physical minutes enter only through the empirical free-flow
and congested speeds, $v^{0}_{kl}$ and $v_{kl}$.

\subsection{Population and fundamentals}

A unit mass $\Lbar$ of workers is allocated across residences
$\{L_i^R\}$ and workplaces $\{L_j^F\}$ with
$\sum_i L_i^R=\sum_j L_j^F=\Lbar$.  Every cell is endowed with a
baseline residential amenity $\ubar_i>0$ and a baseline workplace
productivity $\Abar_j>0$.  Endogenous amenity and productivity load
the local scale of activity through
\begin{equation}
  u_i\;=\;\ubar_i\bigl(L_i^R\bigr)^{\beta},
  \qquad
  A_j\;=\;\Abar_j\bigl(L_j^F\bigr)^{\alpha},
  \qquad \alpha>0,\ \beta<0.
  \label{eq:fundamentals}
\end{equation}
Firms operate under perfect competition with labour as the only
input, so the market-clearing wage at workplace $j$ is
\begin{equation}
  w_j\;=\;A_j\;=\;\Abar_j\bigl(L_j^F\bigr)^{\alpha}.
  \label{eq:wage}
\end{equation}
The composite amenity $u_i$ absorbs the equilibrium effect of
housing-market clearing; goods trade and non-traded amenities are
not modelled separately, consistent with the Allen--Arkolakis
reduced-form parameterisation.

\subsection{Parameters}

The model is closed by four scalar parameters:
\begin{itemize}[leftmargin=2em,itemsep=0.2em]
  \item $\theta>0$, the Fr\'echet shape parameter governing route and
        residence--workplace dispersion;
  \item $\alpha>0$, the production externality (agglomeration)
        elasticity in \eqref{eq:fundamentals};
  \item $\beta<0$, the residential externality (congestion)
        elasticity in \eqref{eq:fundamentals};
  \item $\lambda\ge 0$, the congestion elasticity of speed with
        respect to traffic per lane in \eqref{eq:iceberg}.
\end{itemize}
$(\theta,\alpha,\beta)$ are externally calibrated; $\lambda$ is
internally estimated from a directed speed--flow regression
(\S\ref{ssec:lambda_est}).

%====================================================================
\section{Equilibrium objects}\label{sec:objects}
%====================================================================

This section assembles the equations used in the equilibrium
definition.  Workers choose a residence, a workplace, and a route;
firms hire under perfect competition; the network technology maps
directed link traffic into directed link costs.  Each object is
introduced in turn.

\paragraph{Indirect utility.}
A worker $\nu$ chooses a residence $i$, a workplace $j$, and a route
$r=(r_0=i,r_1,\dots,r_K=j)$ with $(r_{m-1},r_m)\in\E$:
\begin{equation}
  V_{ij,r}(\nu)
  \;=\;
  \frac{u_i\,w_j}{\prod_{m=1}^{K(r)}t_{r_{m-1},r_m}}\,
  \varepsilon_{ij,r}(\nu),
  \qquad
  \varepsilon_{ij,r}\overset{\mathrm{iid}}{\sim}\mathrm{Fr\acute echet}(\theta).
  \label{eq:utility}
\end{equation}

\paragraph{Bilateral commuting cost index.}
Sum over feasible routes:
\begin{equation}
  \tijk^{-\theta}
  \;\equiv\;
  \sum_{r\in\R_{ij}}\prod_{m=1}^{K(r)}t_{r_{m-1},r_m}^{-\theta}.
  \label{eq:tau_def}
\end{equation}
Collect link costs in the directed weighted adjacency matrix
$A=[a_{kl}]$ with $a_{kl}=\tkl^{-\theta}\mathbf 1\{(k,l)\in\E\}$.
Under the spectral-radius assumption below, the series converges and
\begin{equation}
  \tijk^{-\theta}\;=\;\bigl[(I-A)^{-1}\bigr]_{ij}.
  \label{eq:tau_series}
\end{equation}
The bilateral cost index is generically asymmetric,
$\tijk\ne\tau_{ji}$.

\paragraph{Aggregate welfare scale.}
Define
\begin{equation}
  \Wbar\;\equiv\;
  \Bigl[\sum_{(i,j)\in\N^2}\tijk^{-\theta}\,u_i^{\theta}\,w_j^{\theta}\Bigr]^{1/\theta}.
  \label{eq:Wbar_def}
\end{equation}
This is the deterministic welfare index of a resident in the city.
% Ex-ante expected welfare is $\mathbb W=\Gamma(1-1/\theta)\Wbar$.

\paragraph{Gravity equation.}
The bilateral commuting flow is
\begin{equation}
  L_{ij}\;=\;\Lbar\cdot
  \frac{\tijk^{-\theta}\,u_i^{\theta}\,w_j^{\theta}}{\Wbar^{\theta}},
  \qquad
  L_i^R=\sum_j L_{ij},\quad L_j^F=\sum_i L_{ij}.
  \label{eq:gravity}
\end{equation}

\paragraph{Market-access indices.}
Define
\begin{equation}
  \Pi_i^{-\theta}\;\equiv\;\sum_{j}\tijk^{-\theta}\,L_j^F\,P_j^{\theta},
  \qquad
  P_j^{-\theta}\;\equiv\;\sum_{i}\tijk^{-\theta}\,L_i^R\,\Pi_i^{\theta}.
  \label{eq:MA_def}
\end{equation}
Substituting into \eqref{eq:gravity} delivers the equivalent
market-access form of gravity,
\begin{equation}
  L_{ij}\;=\;\tijk^{-\theta}\cdot
  \frac{L_i^R}{\Pi_i^{-\theta}}\cdot
  \frac{L_j^F}{P_j^{-\theta}}.
  \label{eq:gravity_MA}
\end{equation}
Because $\tijk$ is directional, $\Pi_i\ne P_i$ in general.

\paragraph{Link traffic.}
Directed link traffic $\Xi_{kl}$ is the expected number of traversals
of edge $(k,l)$ summed over all commuting flows.  The route system
implies the reduced-form identity
\begin{equation}
  \Xi_{kl}\;=\;\tkl^{-\theta}\,P_k^{-\theta}\,\Pi_l^{-\theta}.
  \label{eq:xi_gravity}
\end{equation}

\paragraph{Congestion technology.}
Given directed link traffic $\Xi_{kl}$, the congestion technology
maps traffic per lane into the directed iceberg link cost
\begin{equation}
  \tkl
  \;=\;
  \tbar_{kl}^{1/\theta}\,
  \bigl(\Xi_{kl}/n_{kl}\bigr)^{\lambda/\theta}.
  \label{eq:cong}
\end{equation}
This is a conditional technology equation rather than an explicit
reduced-form solution, because $\Xi_{kl}$ is itself determined in
equilibrium by \eqref{eq:xi_gravity}.  Substituting
\eqref{eq:xi_gravity} into \eqref{eq:cong} and solving
for $\tkl$ delivers the reduced-form expression
\begin{equation}
  \tkl
  \;=\;
  \tbar_{kl}^{1/[\theta(1+\lambda)]}\cdot
  n_{kl}^{-\lambda/[\theta(1+\lambda)]}\cdot
  \bigl(P_k\Pi_l\bigr)^{-\lambda/(1+\lambda)}.
  \label{eq:tkl_reduced}
\end{equation}

\paragraph{Tidal-lane channel.}
A tidal-lane policy reallocates lanes on a target set
$\E^{\star}\subseteq\E$ from off-peak to peak,
\begin{equation}
  n_{kl}^{\cf}=n_{kl}+\Delta_{kl},\qquad
  n_{lk}^{\cf}=n_{lk}-\Delta_{kl},\qquad
  (k,l)\in\E^{\star},
\end{equation}
preserving total corridor capacity $n_{kl}+n_{lk}$.  Because each
direction enters \eqref{eq:cong} separately, $\tkl$ falls and
$t_{lk}$ rises strictly on treated corridors; the induced changes in
$\{\tijk,L_{ij},L^R,L^F\}$ are asymmetric across $(i,j)$.  Within an
undirected model the policy is vacuous.

\paragraph{Population shares.}
Let $l_i^R\equiv L_i^R/\Lbar$ and $l_j^F\equiv L_j^F/\Lbar$.
Combining \eqref{eq:fundamentals}, \eqref{eq:wage},
\eqref{eq:gravity}, and the marginal-population identities yields
\begin{align}
  (l_i^R)^{1-\beta\theta}
  &\;=\;\chi\sum_{j}\tijk^{-\theta}\,\ubar_i^{\theta}\,\Abar_j^{\theta}\,
       (l_j^F)^{\alpha\theta},
  \label{eq:eq_R}\\
  (l_j^F)^{1-\alpha\theta}
  &\;=\;\chi\sum_{i}\tijk^{-\theta}\,\ubar_i^{\theta}\,\Abar_j^{\theta}\,
       (l_i^R)^{\beta\theta},
  \label{eq:eq_F}
\end{align}
with the scalar $\chi\equiv\Lbar^{\theta(\alpha+\beta)}/\Wbar^{\theta}$
pinned down by the normalisation
$\prod_i\ubar_i=\prod_j\Abar_j=1$.

%====================================================================
\section{Definition of Equilibrium}\label{sec:eq}
%====================================================================

\begin{definition}[Equilibrium]\label{def:eq}
Given primitives $\{(\ubar_i,\Abar_j)\}_{i,j\in\N}$ and
$\{(\ell_{kl},v^{0}_{kl},n_{kl})\}_{(k,l)\in\E}$, parameters
$(\theta,\alpha,\beta,\lambda)$, and total population $\Lbar$, an
equilibrium is the tuple
\[
  \mathcal Q
  \;\equiv\;
  \bigl(
    \{L_i^R\},\,\{L_j^F\},\,\{L_{ij}\},\,
    \{\tkl\},\,\{\tijk\},\,\{\Xi_{kl}\},\,
    \{\Pi_i\},\,\{P_j\},\,\Wbar
  \bigr)
\]
that satisfies, jointly,
\begin{enumerate}[label=(\textbf{E\arabic*}),leftmargin=3em,itemsep=0.25em]
\item \textbf{Bilateral cost index.} \eqref{eq:tau_series}:
  $\tijk^{-\theta}=[(I-A)^{-1}]_{ij}$, with
  $A_{kl}=\tkl^{-\theta}\mathbf 1\{(k,l)\in\E\}$.
\item \textbf{Gravity flows.} \eqref{eq:gravity}:
  $L_{ij}=\Lbar\,\tijk^{-\theta}u_i^{\theta}w_j^{\theta}/\Wbar^{\theta}$,
  with $u_i,w_j$ from \eqref{eq:fundamentals}--\eqref{eq:wage} and
  $\Wbar$ from \eqref{eq:Wbar_def}.
\item \textbf{Market clearing.} $L_i^R=\sum_j L_{ij}$,
  $L_j^F=\sum_i L_{ij}$, $\sum_i L_i^R=\sum_j L_j^F=\Lbar$.
\item \textbf{Market-access indices.} \eqref{eq:MA_def} for
  $(\Pi_i,P_j)$.
\item \textbf{Link traffic.} \eqref{eq:xi_gravity}:
  $\Xi_{kl}=\tkl^{-\theta}P_k^{-\theta}\Pi_l^{-\theta}$.
\item \textbf{Congestion technology.} \eqref{eq:cong}:
  $\tkl=\tbar_{kl}^{1/\theta}(\Xi_{kl}/n_{kl})^{\lambda/\theta}$.
\item \textbf{Spectral-radius condition.} 
\end{enumerate}
% Equivalently, equilibrium share allocations $(l^R,l^F)$ solve the fixed-point system
\eqref{eq:eq_R}--\eqref{eq:eq_F} jointly with
\eqref{eq:tkl_reduced}.
\end{definition}


\paragraph{Welfare representation.}
Two equivalent forms hold at any equilibrium:
\begin{equation}
  \Wbar^{\theta}
  \;=\;\sum_i u_i^{\theta}\Pi_i^{-\theta}
  \;=\;\sum_j w_j^{\theta}P_j^{-\theta}.
  \label{eq:welfare_MA}
\end{equation}
Equation \eqref{eq:welfare_MA} is the welfare object the model
identifies; counterfactual welfare changes are constructed from
exact-hat algebra (\S\ref{ssec:hat}) and reported as the time-equivalent
variation $\hat{\mathbb W}$.  No separate housing-market or goods-market
welfare component is computed.

%====================================================================
\section{Variables, parameters, and data sources}\label{sec:inventory}
%====================================================================

Each object in the empirical implementation is assigned to one of
four roles: data primitive, calibrated or estimated parameter,
inverted fundamental, or endogenous equilibrium variable.  The
distinction fixes what is held fixed across counterfactuals and what
is recomputed after a lane reallocation.

\subsection{Classification convention}\label{ssec:classification}

\begin{itemize}[leftmargin=2em,itemsep=0.2em]
\item \textbf{Data primitive (P)}: constructed from network, speed,
  and commuting data, then held fixed unless the counterfactual
  directly perturbs it.  The lane count $n_{kl}$ is both a primitive
  and the policy lever.
\item \textbf{Parameter (X/I)}: externally calibrated
  $(\theta,\alpha,\beta)$ or internally estimated $(\lambda)$.
\item \textbf{Inverted fundamental (V)}: baseline amenities and
  productivities $(\ubar_i,\Abar_j)$ recovered so that the model
  rationalises the observed residence and workplace marginals.
\item \textbf{Endogenous variable (Q)}: equilibrium objects
  recomputed from Definition~\ref{def:eq}, including commuting
  flows, link traffic, market-access terms, link costs, and welfare.
\end{itemize}

\subsection{Source and construction}\label{ssec:inventory}

Table~\ref{tab:inventory} records the empirical source chain.  Road
geometry, directed lane proxies, and high-frequency speeds are first
matched to the baseline directed centerline network.  These matched
centerline objects are then projected onto the active grid graph to
construct directed edge length, lanes, free-flow speed, and observed
peak speed.  The commuting side starts from the gridded origin--
destination matrix, which is aggregated to the same spatial cells.
The model then estimates or calibrates parameters, inverts the
baseline fundamentals, and recomputes all endogenous variables for
each counterfactual lane allocation.

\begin{small}
\begin{longtable}{p{2.0cm}p{4.2cm}cp{6.5cm}}
\caption{Model variables, parameters, and data sources.}
\label{tab:inventory}\\
\toprule
Symbol & Object & Type & Source / construction \\
\midrule
\endfirsthead
\multicolumn{4}{l}{\emph{(continued)}}\\
\toprule
Symbol & Object & Type & Source / construction \\
\midrule
\endhead
\midrule
\multicolumn{4}{r}{\emph{(continued on next page)}}\\
\endfoot
\bottomrule
\endlastfoot
\multicolumn{4}{l}{\textbf{A. Spatial primitives}}\\
$i,j\in\N$ & Spatial cell (node) & P & Baseline square grid; H3 hex and network-adapted Voronoi grids are robustness tessellations.\\
$(k,l)\in\E$ & Directed grid edge & P & Directed adjacency in the largest connected grid graph, built from the grid travel-cost map of \S\ref{ssec:rawgrid_step}.\\
$\ell_{kl}$ & Directed grid length & P & On-network path length between adjacent grid cells, with access distances added at the endpoints.\\
$n_{kl}$ & Directed lane count & P / $\Delta$ & Length-weighted lane proxy from matched centerlines; this object is perturbed in tidal-lane counterfactuals.\\
$v^{0}_{kl}$ & Free-flow directed speed & P & Nighttime matched-centerline speed, projected to the grid edge.\\
$v^{\obs}_{kl}$ & Observed congested speed & P & Peak-period matched-centerline speed, projected to the grid edge.\\
$\tbar_{kl}$ & Free-flow time & P & $\ell_{kl}/v^{0}_{kl}$.\\
$\tkl^{\obs}$ & Observed iceberg link cost & P & Travel time from $\ell_{kl}$ and $v_{kl}^{\obs}$, converted to iceberg form by \eqref{eq:iceberg}.\\
\midrule
\multicolumn{4}{l}{\textbf{B. Population primitives}}\\
$L_{ij}^{\obs}$ & Bilateral commuting flow & P & Gridded OD matrix aggregated by point-in-polygon to the active tessellation.\\
$L_i^{R,\obs}$ & Residence marginal & P & Row sum $\sum_j L_{ij}^{\obs}$.\\
$L_j^{F,\obs}$ & Workplace marginal & P & Column sum $\sum_i L_{ij}^{\obs}$.\\
$\Lbar$ & Total commuters & P & $\sum_i L_i^{R,\obs}=\sum_j L_j^{F,\obs}$ on the largest weakly connected component of $(\N,\E)$.\\
\midrule
\multicolumn{4}{l}{\textbf{C. Parameters}}\\
$\theta$ & Fr\'echet shape & X & Calibrated from \citet{Ahlfeldt2015}; the internal gravity check is reported in \S\ref{ssec:theta_check}.\\
$\alpha$ & Production externality elasticity & X & Calibrated from the urban productivity literature.\\
$\beta$ & Residential externality elasticity & X & Calibrated from the AA-class spatial-equilibrium literature.\\
$\lambda$ & Congestion elasticity & I & Estimated from the directed speed--flow regression \eqref{eq:lambda_regression}.\\
\midrule
\multicolumn{4}{l}{\textbf{D. Inverted fundamentals}}\\
$\ubar_i$ & Baseline residential amenity & V & Inverted from baseline marginal populations and bilateral commuting costs through \eqref{eq:invert_R}--\eqref{eq:invert_F}.\\
$\Abar_j$ & Baseline workplace productivity & V & Inverted jointly with $\ubar_i$, under the same normalisation.\\
\midrule
\multicolumn{4}{l}{\textbf{E. Endogenous equilibrium variables}}\\
$L_i^R, L_j^F$ & Residence / workplace populations & Q & Solved from \eqref{eq:eq_R}--\eqref{eq:eq_F}; matched to observed marginals at baseline by Assumption~\ref{ass:obs_eq}.\\
$L_{ij}$ & Bilateral commuting flow & Q & Determined by the gravity equation \eqref{eq:gravity}.\\
$u_i,w_j$ & Composite amenity / wage & Q & Determined by \eqref{eq:fundamentals}--\eqref{eq:wage}.\\
$\tijk$ & Bilateral cost index & Q & Computed from directed link costs using \eqref{eq:tau_series}.\\
$\Xi_{kl}$ & Directed link traffic & Q & Computed from \eqref{eq:xi_gravity}; the empirical baseline analogue is constructed in \S\ref{ssec:link_traffic_obs}.\\
$\Pi_i, P_j$ & Resident / firm market access & Q & Linear system \eqref{eq:MA_def}.\\
$\tkl$ & Directed iceberg link cost & Q & Updated through the congestion relation \eqref{eq:cong}.\\
$\Wbar$ & Aggregate welfare scale & Q & Computed from \eqref{eq:Wbar_def} or \eqref{eq:welfare_MA}; counterfactual changes are reported in \S\ref{ssec:hat}.\\
\end{longtable}
\end{small}

%====================================================================
\section{Estimation and Algorithms}\label{sec:est}
%====================================================================

The empirical pipeline executes in five steps:
\begin{enumerate}[label=\textbf{Step \arabic*.},leftmargin=3.5em,itemsep=0.2em]
\item Construct the directed primitives
  $(\ell_{kl},v^{0}_{kl},v_{kl}^{\obs},n_{kl})$ and the population
  primitives $(L_{ij}^{\obs},L_i^{R,\obs},L_j^{F,\obs},\Lbar)$
  (\S\ref{ssec:data_construction}).
\item Calibrate $(\theta,\alpha,\beta)$ externally
  (\S\ref{ssec:calibration}).
\item Recover or calibrate the congestion elasticity from speed--flow
  variation
  (\S\ref{ssec:lambda_est}).
\item Invert $(\ubar_i,\Abar_j)$ from the equilibrium system at the
  observed allocation (\S\ref{ssec:invert}).
\item Solve counterfactuals with the AA exact-hat algorithm
  (\S\ref{ssec:hat}).
\end{enumerate}
Each step is described below as inputs $\to$ algorithm $\to$ outputs.

\subsection{Step 1: Constructing data inputs}\label{ssec:data_construction}

The model takes the following directly constructed objects as inputs:
directed grid-edge length $\ell_{kl}$, free-flow speed $v_{kl}^{0}$,
observed peak-period speed $v_{kl}^{\obs}$, directed lane capacity
$n_{kl}$, bilateral commuting flows $L_{ij}^{\obs}$, residence
marginals $L_i^{R,\obs}$, workplace marginals $L_j^{F,\obs}$, and
total commuter mass $\Lbar$.  The construction of these variables
from raw road, speed, lane, and OD data is documented in the data
construction notes.  This section treats them as primitives for the
estimation and counterfactual steps.  The active directed grid graph
used for these objects is the grid travel-cost map referenced below.
\label{ssec:rawgrid_step}

\subsection{Step 2: External calibration of $(\theta,\alpha,\beta)$}
\label{ssec:calibration}

The baseline calibration follows the Seattle urban application in
\citet{AllenArkolakis2022}, which adopts the quantitative urban
parameters from \citet{Ahlfeldt2015}:
\begin{itemize}[leftmargin=2em,itemsep=0.2em]
	\item $\theta=6.83$, the workplace--residence preference dispersion
	parameter.
	\item $\alpha=-0.12$, the workplace externality parameter used in
	the AA Seattle calibration.
	\item $\beta=-0.10$, the residence externality parameter used in
	the AA Seattle calibration.
\end{itemize}
These values match the current replication code.  Alternative
parameter values are left for robustness exercises rather than used
in the baseline replication.



\subsection{Step 3: Congestion elasticity from speed--flow variation}
\label{ssec:lambda_est}

\paragraph{Theoretical target.}
The congestion technology links link speed to traffic per lane.  In
the notation of Section~\ref{sec:setup},
$v_{kl}=v^{0}_{kl}(n_{kl}/\Xi_{kl})^{\lambda}$, so the estimating
equation is
\begin{equation}
  \ln v_{kl}^{\obs}
  \;=\;
  \ln v_{kl}^{0}
  -\lambda\,\ln\!\bigl(\hat\Xi_{kl}/n_{kl}\bigr)
  +\eta_{kl}
  +\varepsilon_{kl}.
  \label{eq:lambda_regression}
\end{equation}
The clean specification is a panel version of \eqref{eq:lambda_regression}
with directed-link fixed effects and time fixed effects:
\begin{equation}
  \ln v_{e t}
  \;=\;
  \mu_e+\rho_t-\lambda\ln(\Xi_{e t}/n_{e t})+\varepsilon_{e t},
  \qquad e=(k,l).
  \label{eq:lambda_panel}
\end{equation}
Because traffic is chosen in response to congestion, a preferred
specification instruments traffic per lane with shocks that move
traffic or capacity but are orthogonal to contemporaneous link-speed
shocks.  Plausible data sources should include traffic counts and IV
measures.  These data would identify $\lambda$ directly.

\paragraph{Application with current data.}
The present data contain speeds, free-flow speeds, lane proxies, and
commuting OD flows, but not measured directed traffic.  A structural
estimate of $\lambda$ is therefore not available from the current
data alone.\label{ssec:link_traffic_obs}  The implementation uses
two reduced-form substitutes.
First, the AA replication path imports the paper's congestion
coefficient $\delta_1=0.488$ and converts it to the iceberg-cost
parameter $\lambda=\delta_1/\theta$ used by the exact-hat solver.
Second, the diagnostic Python routine
\texttt{estimate\_lambda\_cross\_section} regresses
$\log(t_{kl}^{\obs}/t_{kl}^{0})$ on
$\log(\hat\Xi_{kl}^{\obs}/n_{kl})$ after removing lane-bin means.
The latter is a weak calibration check because
$\hat\Xi_{kl}^{\obs}$ is model-assigned rather than observed.  It is
not treated as the baseline identifying estimate.

\subsection{Step 4: Inversion of $(\ubar_i,\Abar_j)$}\label{ssec:invert}

\begin{assumption}[Observed allocation is an equilibrium]
\label{ass:obs_eq}
$(L_i^{R,\obs},L_j^{F,\obs})$ equals $(L_i^R,L_j^F)$ in
Definition~\ref{def:eq} under primitives
$(\ell_{kl},v_{kl}^{0},n_{kl})$, parameters
$(\theta,\alpha,\beta,\lambda)$, and some baseline fundamentals
$(\ubar_i,\Abar_j)$.  The bilateral cost index used in the inversion
is $\tijk^{\obs}=[(I-A^{\obs})^{-1}]_{ij}^{-1/\theta}$ with
$A^{\obs}_{kl}=(\tkl^{\obs})^{-\theta}$.
\end{assumption}

Under Assumption~\ref{ass:obs_eq}, the equilibrium system
\eqref{eq:eq_R}--\eqref{eq:eq_F} evaluated at the data is
\begin{align}
  (l_i^{R,\obs})^{1-\beta\theta}
  &\;=\;\chi\sum_{j}(\tijk^{\obs})^{-\theta}\,\ubar_i^{\theta}\,\Abar_j^{\theta}\,
       (l_j^{F,\obs})^{\alpha\theta},
  \label{eq:invert_R}\\
  (l_j^{F,\obs})^{1-\alpha\theta}
  &\;=\;\chi\sum_{i}(\tijk^{\obs})^{-\theta}\,\ubar_i^{\theta}\,\Abar_j^{\theta}\,
       (l_i^{R,\obs})^{\beta\theta}.
  \label{eq:invert_F}
\end{align}

Equations~\eqref{eq:invert_R}--\eqref{eq:invert_F} are solved by
alternating updates of $\ubar_i^\theta$ and $\Abar_j^\theta$ with a
normalisation on the two geometric means.  The output is the pair
$(\ubar_i^{\star},\Abar_j^{\star})$, which is held fixed in
counterfactual analysis.  This step is an inversion of baseline
fundamentals, not a separate estimation of amenities or productivity.

\subsection{Step 5: Counterfactual exact-hat solver}\label{ssec:hat}

Counterfactuals are implemented as shocks to bilateral iceberg
commuting costs.  The solver takes as input a matrix
$\hat{\bar t}_{ij}$, where each element is the ratio of the
counterfactual cost to the baseline cost.  In the one-edge validation
exercise, a treated directed edge receives
$\hat{\bar t}_{ij}=0.99$, with all other entries equal to one.  In
the tidal-lane exercise, the treated directions receive
\[
\hat{\bar t}_{ij}
=
\frac{\tau_{ij}^{\mathrm{new}}}{\tau_{ij}^{\obs}},
\]
where $\tau_{ij}^{\mathrm{new}}$ is computed from the BPR
lane-reallocation rule using baseline traffic and counterfactual lane
capacity.

Given $\hat{\bar t}_{ij}$, the AA exact-hat solver recomputes the
counterfactual equilibrium in ratios rather than levels.  Baseline
traffic, residence masses, workplace masses, and OD flows are read
from the AA-format input files; fundamentals are held fixed.  The
solver then iterates on the residence and workplace population hats
$(\hat l_R,\hat l_F)$ and the equilibrium scale term $\hat\chi$
until the counterfactual system converges.

The reported outputs are the welfare-related scale change
$\hat\chi$, the population changes $(\hat l_R,\hat l_F)$, and the
model-implied changes in travel costs and traffic.  In this
normalisation, $\hat\chi<1$ corresponds to a welfare gain.


\subsection{Counterfactual scenarios}\label{ssec:scenarios}

\begin{description}[leftmargin=2.5em,itemsep=0.25em]
\item[CF1.] \emph{Symmetric vs.\ directional commuting costs.} \\ 
  Replace $\tkl^{\obs}$ by the corridor-level average
  $\tfrac12(\tkl^{\obs}+t_{lk}^{\obs})$, holding lane capacities and
  fundamentals fixed.  Isolates the welfare contribution of observed
  directionality.
\item[CF2.] \emph{Tidal-lane reallocation.}  \\
  Apply
  $n_{kl}^{\cf}=n_{kl}+\Delta$ and
  $n_{lk}^{\cf}=n_{lk}-\Delta$ for selected bidirectional pairs,
  where $(k,l)$ is the more congested peak-period direction.  The
  treated set is obtained from a directed speed-asymmetry screen, and
  the lane transfer preserves total corridor capacity
  $n_{kl}+n_{lk}$. (TO DISCUSS: TRAGETS?)
\item[CF3.] \emph{Uniform lane expansion.} Multiply $n_{kl}$ by a
  factor in $\{1.10,1.25,1.50\}$ uniformly across links.  Benchmarks
  CF2 against undifferentiated capacity addition.
\end{description}
Across CF1--CF3, the reported outcomes are the welfare-related
scale change $\hat\chi$, residence and workplace population changes
$(\hat l_R,\hat l_F)$, and model-implied changes in travel costs and
traffic.  Robustness exercises vary the peak period, grid system,
and calibrated elasticities.

\subsection{Internal robustness check on $\theta$}\label{ssec:theta_check}

For each $\theta\in\{3,4,\dots,12\}$ the procedure recomputes
$\tijk^{\obs}(\theta)$ via \eqref{eq:tau_series} and runs the
two-way fixed-effects gravity regression
\begin{equation}
  \ln L_{ij}^{\obs}\;=\;-\theta\,\ln\tijk(\theta)+\mu_i+\nu_j+\varepsilon_{ij}.
  \label{eq:gravity_log}
\end{equation}
The internal-consistency estimator selects the $\theta$ at which the
regression slope equals the imposed $\theta$.  The result is
reported alongside the external calibration; the external calibration
remains the baseline because \eqref{eq:gravity_log} relies on the
same network costs that are themselves equilibrium objects.

%====================================================================
\begin{thebibliography}{10}
%====================================================================
\bibitem[Ahlfeldt et al.(2015)]{Ahlfeldt2015}
Ahlfeldt, G., Redding, S. J., Sturm, D. M., Wolf, N. (2015). The
economics of density: Evidence from the Berlin Wall.
\emph{Econometrica}, 83(6), 2127--2189.

\bibitem[Allen and Arkolakis(2022)]{AllenArkolakis2022}
Allen, T., Arkolakis, C. (2022). The welfare effects of transportation
infrastructure improvements. \emph{Review of Economic Studies},
89(6), 2911--2957.

\end{thebibliography}

\end{document}
